%\documentclass[overheadsbeamer.tex]{subfiles}
\documentclass[handoutsbeamer.tex]{subfiles}

\begin{document}

\thispagestyle{empty}

\ona{ \setcounter{chapter}{1} } % ONE LESS
\lecture{The Exponential Quantum Advantage Hypothesis}{Lect02}
\chapter{The Exponential Quantum Advantage Hypothesis}\label{C.EQA}

\ona{This \chap{} examines the claim that motivated much of quantum computing for chemistry, and by extension the variational methods of this course: that quantum computers offer an \emph{exponential} advantage for the ground-state problem of molecules. We follow the critical evaluation of \citet{lee2023evaluating} and place it next to the complexity-theoretic view of the survey \cite{abbas2024challenges}. The lesson carries over to optimization: worst-case complexity is neither necessary nor sufficient for a practical advantage.}

\section{The hypothesis}

\begin{frame}\ft{What is claimed}
\ona{The ground-state energy of the electronic Hamiltonian is the central quantity of quantum chemistry (\Chap~\chref{C.Motivation}). Exact classical methods (full configuration interaction) scale exponentially in the number of electrons and orbitals. Quantum phase estimation (QPE), given a trial state $\ket{\phi}$ with overlap $|\braket{\phi|\psi_0}|^2 = \gamma^2$ with the true ground state, returns the ground-state energy to precision $\epsilon$ with $\mathcal O(1/(\gamma^2\epsilon))$ uses of a Hamiltonian simulation of polynomial cost. This suggested:}
\begin{center}
\emph{The exponential quantum advantage (EQA) hypothesis:} for generic ground-state chemistry problems, quantum algorithms achieve an exponential speedup over the best classical algorithms.
\end{center}
\ona{The hypothesis has two halves, both of which must hold for \emph{typical} molecules of interest: (i) the quantum algorithm runs in polynomial time, and (ii) every classical algorithm requires exponential time.}
\onp{\begin{itemize}\item QPE: cost $\mathcal O(1/(\gamma^2\epsilon))$, $\gamma$ = overlap of the trial state with the ground state. \item Exact classical methods: exponential. \item EQA: exponential speedup for \emph{generic} ground-state chemistry.\end{itemize}}
\end{frame}

\begin{frame}\ft{What is known from complexity theory}
\begin{itemize}\setlength\itemsep{0.3em}
    \item The general local-Hamiltonian ground-state problem is $\mathbf{QMA}$-complete: hard even for quantum computers. Electronic-structure Hamiltonians are a restricted family, but the ground-state problem with an arbitrary external potential remains $\mathbf{QMA}$-hard.
    \item Hence no polynomial-time quantum algorithm can solve \emph{all} instances, and half (i) of EQA can only hold for the instances chemists care about, i.e., it is a statement about \emph{typical} or \emph{structured} instances, not worst cases.
    \item Half (ii) is a statement about the absence of classical algorithms, which complexity theory cannot deliver for typical instances either: classical heuristics such as coupled cluster, density-matrix renormalisation group and quantum Monte Carlo are not exact, but succeed on large classes of molecules.
\end{itemize}
\ona{So the hypothesis is empirical, and \citet{lee2023evaluating} test it empirically, as one tests any hypothesis about heuristics.}
\end{frame}

\section{Evaluating the evidence}

\begin{frame}\ft{Half (i): the overlap problem}
\ona{QPE needs a trial state with non-negligible overlap $\gamma$. For a single molecule the overlap of a mean-field (Hartree--Fock) state is often reasonable. But for the EQA hypothesis one must ask how $\gamma$ scales with system size for a \emph{family} of instances, e.g., a chain of iron--sulphur clusters or a growing lattice of a correlated material.}
\begin{itemize}\setlength\itemsep{0.3em}
    \item For weakly correlated systems, the overlap of a mean-field state decays exponentially with size (the orthogonality catastrophe), but so do the correlation energies per subsystem stabilise, and classical methods succeed.
    \item For strongly correlated systems, \citet{lee2023evaluating} numerically find that overlaps of the best cheap trial states also decay exponentially with size, e.g., in Fe--S clusters and in the Hubbard model away from half filling; more elaborate state-preparation heuristics (adiabatic preparation, variational circuits) may help, but their cost is not understood, and adiabatic preparation itself faces vanishing gaps.
    \item An exponentially small $\gamma$ makes QPE exponentially expensive: half (i) fails unless a polynomial-cost state preparation exists, which is exactly the hard part.
\end{itemize}
\end{frame}

\begin{frame}\ft{What can be guaranteed given the overlap}
\ona{The overlap is also where the positive results live. \citet{zhang2022computing} show that if one \emph{is} given a state with overlap $\eta$ and a Hamiltonian with gap $\gamma$, then not only the ground-state energy \cite{lt21} but any ground-state property $\braket{\psi_0|O|\psi_0}$ can be estimated to accuracy $\epsilon$ with maximal evolution time $\widetilde{\mathcal O}(\gamma^{-1})$ and total evolution time $\widetilde{\mathcal O}(\gamma^{-1}\epsilon^{-2}\eta^{-2})$, using one ancilla; see \Chap~\chref{C.Motivation}. Two remarks:}
\begin{itemize}\setlength\itemsep{0.3em}
    \item The guarantee is conditional on $\eta$, and the cost scales as $\eta^{-2}$: the EQA question is precisely whether $\eta$ stays polynomial along families of instances.
    \item Even with the assumptions, the task is not classically easy: estimating ground-state properties to high precision classically would imply $\mathbf P = \mathbf{BQP}$; and without the assumptions, the decision version is $\mathsf{P^{QMA[\log]}}$-complete, harder than energy estimation. So the \emph{form} of a possible advantage is clear: polynomial cost given a good trial state, on instances where classical heuristics do not converge. Its existence for natural families is what remains unproven.
\end{itemize}
\end{frame}

\begin{frame}\ft{Half (ii): are classical heuristics really exponential?}
\ona{Exponential classical cost is easy to prove for exact methods and hard to prove for heuristics. \citet{lee2023evaluating} instead \emph{measure} the empirical scaling of classical heuristics to chemical accuracy:}
\begin{itemize}\setlength\itemsep{0.3em}
    \item For the molecules and materials typically cited as motivation (transition-metal complexes, Fe--S clusters, cytochrome P450, FeMoco), the cost of coupled-cluster, DMRG and selected-CI methods to reach chemical accuracy grows polynomially in practice, with modest exponents, across the sizes accessible.
    \item In the few settings where classical heuristics do scale exponentially, the quantum algorithm's state-preparation cost appears to scale exponentially too, for the same physical reason: the ground state has little structure that either side can exploit.
    \item The quantitative gap between polynomial classical scaling and polynomial quantum scaling, once constant factors and error correction are included, leaves at most a \emph{polynomial} speedup, and often a disadvantage.
\end{itemize}
\ona{The conclusion of \citet{lee2023evaluating} is that the evidence for generic exponential quantum advantage in ground-state chemistry is lacking; polynomial speedups remain possible and would still be valuable.}
\end{frame}

\begin{frame}\ft{The picture in one slide}
\begin{center}\ona{\small}\onp{\scriptsize}
\begin{tabular}{lcc}\toprule
 & classical heuristics & QPE with cheap trial state \\\midrule
weakly correlated & polynomial, accurate & polynomial \\
strongly correlated, structured & polynomial in practice & polynomial if $\gamma$ large \\
strongly correlated, unstructured & exponential & exponential ($\gamma\to0$) \\\bottomrule
\end{tabular}
\end{center}
\ona{Exponential advantage would require the lower-left to be exponential while the lower-right is polynomial, i.e., instances that are hard for every classical heuristic but admit an easily prepared trial state. No natural family of molecules with this property is known.}
\onp{EQA needs: classical exponential \emph{and} quantum polynomial on the same instances. No natural family is known.}
\end{frame}

\section{Heuristics, formally, and what this means for optimization}

\begin{frame}\ft{Heuristics, formally}
\ona{The same tension appears in optimization, where the survey \cite{abbas2024challenges} makes the point precisely. Heuristic algorithms are algorithms without provable performance guarantees, run-time guarantees, or both. One way to make claims about them precise is distributional: a pair $(L,D)$ of a language and a family of input distributions is in $\mathbf{HeurBPP}$ if a polynomial-time randomized algorithm $\mathcal A$ satisfies, for all $n$ and $\delta > 0$,}
\begin{equation*}
    \mathrm{Pr}_{x \sim D_n}\Big[\mathrm{Pr}\big[\mathcal{A}(x,0^{1/\delta}) = L(x)\big] \ge 2/3\Big] \ge 1-\delta,
\end{equation*}
\ona{i.e., it may err on a small fraction of instances; $\mathbf{HeurBQP}$ is the quantum analogue.}
\begin{itemize}\setlength\itemsep{0.2em}
    \item Super-polynomial separations between the two are known only for contrived distributions, essentially Karp-reduction images of factoring; they have no practical value.
    \item The overlap gap property makes some random constraint satisfaction problems average-case hard for large classes of classical \emph{and} low-depth quantum algorithms (\Chap~\chref{C.Inapprox}); whether deeper quantum algorithms escape is open.
\end{itemize}
\end{frame}

\begin{frame}\ft{What complexity theory does and does not tell us}
\ona{It is widely believed that quantum computers will not offer exponential speedups for $\mathbf{NPO}$-hard problems in the worst case. But in practice one rarely needs the worst case or an exact solution; approximations are acceptable, and the hardness of \emph{approximation} with quantum computers is largely unexplored. The survey's summary applies to chemistry and optimization alike:}
\begin{center}
    \textit{Understanding complexity theory is extremely useful for gauging a possible quantum advantage, but rigorous complexity-theoretic separations are neither necessary nor sufficient when seeking a practical quantum advantage.}
\end{center}
\ona{TSP illustrates the nuance: its decision version is $\mathbf{NP}$-hard and the optimization version $\mathbf{APX}$-hard; Euclidean TSP has a polynomial-time approximation scheme; heuristics solve instances with millions of nodes to near-optimality; and shortest path is in $\mathbf{P}$. In contrast, there exist problems with fewer than 100 variables that are hard classically, such as low-autocorrelation binary sequences.}
\end{frame}

\begin{frame}\ft{Lessons for variational methods}
\begin{enumerate}\setlength\itemsep{0.4em}
    \item \textbf{Claims must be about instance families}, with the classical baseline being the best heuristic, not the exact method. For MaxCut that baseline is Goemans--Williamson plus local search, not brute force (\Chap~\chref{C.Motivation}).
    \item \textbf{State preparation is the crux.} In chemistry it is the overlap $\gamma$; in QAOA it is the initial state and the angles. Warm starts (\Chaps~\chref{C.WSRounded}--\chref{C.WSContinuous}) are the optimization analogue of a good trial state.
    \item \textbf{Benchmark empirically and fairly}: model-independent problems, time-to-solution and cost, distributions over instances, and validation on hardware (\Chap~\chref{C.PerIteration}).
    \item \textbf{Polynomial advantages are worth having}, and small improvements matter in applications such as finance (\Chap~\chref{C.Portfolio}); but they must be measured, not assumed.
\end{enumerate}
\end{frame}

\begin{frame}[allowframebreaks]\ft{References}
\biblio
\end{frame}

\end{document}
